% AUTO-GENERATED by scripts/exp_sync_kernel_real_input_40_arity_3d.py
% Derived consequences for the (3,3,5) tensor with third axis N2 mod 5.
\paragraph{从 $((3,3,5))$ 输出推导出的二级结论（可复现）}
下述量全部由上表的 5 张 $3\times 3$ 切片与最坏扭曲谱半径比直接计算得到。
\paragraph{（i）慢模完全由碰撞 cocycle 主导}
最坏谱半径比与达到最坏的角色索引为：
$$
\frac{\rho_{3,3,5}}{\lambda}\approx 0.950621088873.
$$
$$
(0,0,1),\ (0,0,4)
$$
由此定义主相关长度尺度（指数衰减的时间常数）与半衰期：
$$
\tau_{\mathrm{mix}}:=\frac{1}{-\log(\rho_{3,3,5}/\lambda)}\approx 19.747340579191,\qquad t_{1/2}:=\frac{\log 2}{-\log(\rho_{3,3,5}/\lambda)}\approx 13.687813446023.
$$
\paragraph{（ii）碰撞同余类 $c$ 的常数漂移（总偏置）}
定义 $S_c:=\sum_{a=0}^2\sum_{b=0}^2 C^{(3,3,5)}_{a,b,c}$。数值上
$$
(S_0,S_1,S_2,S_3,S_4)\approx(+1.685330,+0.130189,-0.321067,-0.528302,-0.643249).
$$
\paragraph{（iii）三维耦合在 $c$ 轴上的局部化（切片方差）}
定义每张切片在 $3\times 3$ 网格上的标准差（按 $9$ 个条目平均）：
$$
\bigl(\mathrm{sd}(c=0),\dots,\mathrm{sd}(c=4)\bigr)\approx\bigl(0.347585,0.0591959,0.00515942,0.000787443,0.000137309\bigr).
$$
\paragraph{（iv）$\mathbb{{Z}}/5$ 轴的傅里叶压缩指纹（均值 + 两个复模）}
对每个固定的 $(a,b)$，令 $f_{a,b}(c):=C^{(3,3,5)}_{a,b,c}$（$c\in\ZZ/5$）。其离散傅里叶系数为
$$
\widehat f_{a,b}(j):=\frac15\sum_{c=0}^{4} f_{a,b}(c)\,\omega_5^{-jc},\qquad \omega_5=e^{2\pi i/5}.
$$
从而有恒等分解（对所有 $c$ 精确成立）
$$
f_{a,b}(c)=\widehat f_{a,b}(0)+2\Re\bigl(\widehat f_{a,b}(1)\,\omega_5^{c}\bigr)+2\Re\bigl(\widehat f_{a,b}(2)\,\omega_5^{2c}\bigr).
$$
记 $\overline C_{b,a}:=\widehat f_{a,b}(0)$（行 $b$、列 $a$）与两张复矩阵 $A^{(1)}_{b,a}:=\widehat f_{a,b}(1)$、$A^{(2)}_{b,a}:=\widehat f_{a,b}(2)$。则：
$$
\overline C=
\begin{pmatrix}
\phantom{-}0.06447040 & \phantom{-}0.15692811 & \phantom{-}0.02518478\\
-0.04778699 & \phantom{-}0.01447759 & -0.04317664\\
-0.01940426 & -0.04655149 & -0.03956088
\end{pmatrix}
$$
并且（以条目绝对值的最大值度量）
$$
\max_{a,b}|A^{(1)}_{a,b}|\approx 0.201317996544,\qquad\max_{a,b}|A^{(2)}_{a,b}|\approx 0.200144063931,\qquad\frac{\max|A^{(2)}|}{\max|A^{(1)}|}\approx 0.994168764674.
$$
其中 $A^{(1)}$ 与 $A^{(2)}$ 的具体条目可由同一脚本在 JSON 输出中复核。
